\chapter{Using the stack}
\label{ch:use}

The stack turns one LaTeX source into PDFs in several layouts and a website. It is a thin layer over LaTeX, Pandoc, PDF.js and Zensical. A project needs one file, \texttt{document.toml}, beside its source.

\section{A guide}

\begin{lstlisting}[language=text]
name = "my_guide"
title = "My guide"
author = "Your Name"
layouts = ["manual"]
source = "tex"
chapters = ["ch_one", "ch_two"]
figures = "figures"
shared = ["../shared"]
[web]
site_name = "My guide"
intro = "One sentence for the home page."
\end{lstlisting}

\section{A manuscript}

A manuscript names its main file, and the stack reads the title, abstract, body and bibliography from it.

\begin{lstlisting}[language=text]
name = "my_paper"
layouts = ["own"]
source = "paper"
main = "main.tex"
derive = "main.tex"
\end{lstlisting}

\section{Commands}

\begin{lstlisting}[language=bash]
python -m research_helpers.stack layouts
python -m research_helpers.stack pdf .
python -m research_helpers.stack pdf . ieee
python -m research_helpers.stack web .
python -m research_helpers.stack check .
python -m research_helpers.stack snapshot .
python -m research_helpers.stack deploy . ../site/teaching/my-guide
\end{lstlisting}

The \texttt{layouts} command lists every layout. A working example is in \texttt{examples/demo}, and the personal site template is at \texttt{github.com/AhmadAmmar/academic-site-template}.
